% Designed by Yiming Xu and Xiangjian Zhang in Gainesville, FL.
% All rights reserved.

% Excel to Latex: https://www.tablesgenerator.com/

\documentclass[final]{beamer}

% ====================
% Packages
% ====================

\usepackage[flushleft]{threeparttable}
\usepackage[T1]{fontenc}
\usepackage{lmodern}
\usepackage[size=custom,width=120,height=72,scale=1.0]{beamerposter}
\usetheme{gemini}
\usecolortheme{uchicago}
\usepackage{graphicx}
\usepackage{booktabs}
\usepackage{tikz}
\usepackage{pgfplots}
\pgfplotsset{compat=1.17}
\usepackage{caption}

% ====================
% Lengths
% ====================

% If you have N columns, choose \sepwidth and \colwidth such that
% (N+1)*\sepwidth + N*\colwidth = \paperwidth
\newlength{\sepwidth}
\newlength{\colwidth}
\setlength{\sepwidth}{0.025\paperwidth}
\setlength{\colwidth}{0.3\paperwidth}

\newcommand{\separatorcolumn}{\begin{column}{\sepwidth}\end{column}}

% ====================
% Title
% ====================

\title{The location problem of EV charging stations: research trends, critical gaps, and a future agenda}

\author{Kelly Anran Zheng \inst{1} {\normalsize (anranzheng@ufl.edu)} \and Xiang 'Jacob' Yan \inst{1} {\normalsize (xiangyan@ufl.edu)}}

\institute[shortinst]{\inst{1} Department of Civil and Coastal Engineering, University of Florida }


% ====================
% Footer (optional)
% ====================

\footercontent{
  \href{ }{ } \hfill
  102nd TRB Annual Meeting, Washington D.C., 2023 \hfill
  \href{}{}}
% (can be left out to remove footer)

% ====================
% Logo (optional)
% ====================

% use this to include logos on the left and/or right side of the header:
% \logoright{\includegraphics[height=7cm]{logo1.pdf}}
% \logoleft{\includegraphics[height=7cm]{logo2.pdf}}

% ====================
% Body
% ====================

\begin{document}
\addtobeamertemplate{headline}{}
{
    \begin{tikzpicture}[remember picture,overlay]
      \node [anchor=north west, inner sep=3cm] at ([xshift=1.2cm,yshift=-3cm]current page.north west)
      {\includegraphics[height=2.5cm]{logos/UFTransportation_Institute.png}}; % also try shield-white.eps
      \node [anchor=north east, inner sep=3cm] at ([xshift=-1cm,yshift=-2.6cm]current page.north east)
      {\includegraphics[height=2.8cm]{logos/logo.png}};
    \end{tikzpicture}
}

\begin{frame}[t]
\begin{columns}[t]
\separatorcolumn

\begin{column}{\colwidth}

  \begin{block}{Background}
  \begin{itemize}
      \item \textcolor[RGB]{7,84,156}{\textbf{Eletric vehicles (EVs)}} are increasingly adopted by consumers around the world, creating an urgent need for the development of \textcolor[RGB]{7,84,156}{\textbf{EV charging infrastructure}} in a variety of contexts. 
      \item Much research work have been undertaken recently regarding \textcolor[RGB]{7,84,156}{\textbf{locating and sizing EV charging stations (EVCS)}}. However, few reviews provided an in-depth analysis of \textcolor[RGB]{7,84,156}{\textbf{key gaps}} in existing research, such as \textcolor[RGB]{7,84,156}{\textbf{data gaps, limitations in approaches and lack of consideration for key issues}}.
  \end{itemize}

  \end{block}

    \begin{block}{Research Objectives}
    \begin{itemize}
      \item Identify \textcolor[RGB]{7,84,156}{\textbf{datasets}} commonly used in current research and the data gaps.
      \item Identify \textcolor[RGB]{7,84,156}{\textbf{critical issues}} related to EV charging infrastructure largely ignored or addressed inadequately by the existing literature.
      \item Review \textcolor[RGB]{7,84,156}{\textbf{existing modeling and analytical approaches}} and discuss their limitations.
      \item Review existing empirical studies to explore if some important \textcolor[RGB]{7,84,156}{\textbf{study contexts }}(e.g., rural areas and underserved communities) have received less attention in existing studies. 
      \item Propose a \textcolor[RGB]{7,84,156}{\textbf{future research agenda}} and \textcolor[RGB]{7,84,156}{\textbf{policy recommendations}} for EVCS development.
    \end{itemize}
  \end{block}

  \begin{block}{Overview of research trends and topics}
We first examine the recent research trends on EVCS with a \textcolor[RGB]{7,84,156}{\textbf{bibliographic analysis and a co-citation analysis}}. We find a growing research interest on four research themes, including \textcolor[RGB]{7,84,156}{\textbf{economics, energy sources, infrastructure planning, and network design}}.

\begin{figure}
  \centering
  \includegraphics[scale=1.2]{logos/paper num.png}
  \hspace{0.01in}
  \includegraphics[scale=1.2]{logos/by country.png}
  \caption{(a). Publications by year since 2000    (b). The top 10 countries with the most publications
}
\end{figure}

With a focus on issues concerning the fields of \textcolor[RGB]{7,84,156}{\textbf{transportation engineering and planning}}, a total of \textcolor[RGB]{7,84,156}{\textbf{152
publications}} were reviewed finally and were classified into \textcolor[RGB]{7,84,156}{\textbf{four categories}}:

1. charging site selection and network optimization

2. charging behavior analysis and modeling

3. charging demand prediction

4. equity in spatial access

  \end{block}

\end{column}

\separatorcolumn

\begin{column}{\colwidth}


\begin{figure}
  \centering
  \includegraphics[scale=1.5]{logos/cluster.png}
  \captionsetup{font={small,bf,stretch=1.25}}
  \caption{Network relationship of most frequent keywords}
\end{figure}


  \begin{block}{Analysis of current research}

\begin{table}[!ht]\tiny%
\newcommand{\tabincell}[1]{}
\caption{Commonly used datasets in EVCS research and major limitations}
\resizebox{.85\columnwidth}{!}{
\centering
\begin{tabular}{lp{10cm}}
\hline
Types of datasets  & Examples         \\ \hline
Infrastructure& Arrival \& departure time, travel duration, travel distance, travel path, origin and destination \\ 
Mobility & Transit routes and stations, parking lots, current \& planned charging stations \\
People & Household income, vehicle purchases \& ownership, race, charging demand \\
Land use & Public facilities POI, topology proximity to facilities, natural disaster risk \\
\hline
\end{tabular}}
\end{table}
\begin{figure}
  \centering
  \includegraphics[scale=1.2]{logos/map.png}
  \caption{Geographical distribution of case study regions in our reviewed work}
\end{figure}

  \end{block}



\end{column}

\separatorcolumn

\begin{column}{\colwidth}
\begin{figure}
  \centering
  \includegraphics[scale=1.28]{logos/trb.drawio.png}
    \captionsetup{font={small,bf,stretch=1.2}}
  \caption{Overview of research objectives and analytical approaches in existing EVCS research}
\end{figure}

    \begin{alertblock}{Key research gaps}
    \begin{itemize}
      \item Due to \textcolor[RGB]{250,70,22}{\textbf{data unavailability}}, \textcolor[RGB]{250,70,22}{\textbf{\textcolor[RGB]{250,70,22}{\textbf{uncertainty characteristics}}}} are usually neglected in location optimization of EVCS. {\textbf{\textcolor[RGB]{250,70,22}{\textbf{Time fluctuations}}}} are less examined in EV charging demand. 
      \item There is still lack of \textcolor[RGB]{250,70,22}{\textbf{high dimension}} with input features for ML model training to plan the allocation of EVCS, e.g. high-resolution driving profiles, trafffc and weather conditions.
      \item Studies on the charging needs of \textcolor[RGB]{250,70,22}{\textbf{disadvantaged population groups}} are limited. 
      \item Limited work has examined how to facilitate \textcolor[RGB]{250,70,22}{\textbf{shared micromobility}} use by deploying shared micro-mobility charging station across different geographic and socioeconomic contexts.

    \end{itemize}
  \end{alertblock}





  \begin{block}{Future research and policy directions}
1. Collect and analyze data with higher spatiotemporal granularity. 

2. Integrate uncertainty factors and temporal variations into EVCS analysis and modeling.

3. Strengthen the linkages between EV charging and the following: future mobility, land use, and social equity. 




  \end{block}


\end{column}

\separatorcolumn
\end{columns}
\end{frame}

\end{document}
